\documentclass{article}
\usepackage[top=2.3cm, bottom=2.3cm, left=2.1cm, right=2.1cm]{geometry}
\usepackage{amsmath,amssymb,amsthm,mathtools}
\usepackage{bbm}
\usepackage{booktabs}
\usepackage{enumitem}
\usepackage{hyperref}
\usepackage{natbib}
\usepackage{iftex}
\usepackage[default,regular,black]{sourceserifpro}
\ifPDFTeX
  \usepackage[T1]{fontenc}
\fi
\hypersetup{colorlinks,linkcolor={blue},citecolor={blue},urlcolor={red}}

% ====== Notation macros ======
\newcommand{\Align}{\mathrm{Align}}
\newcommand{\FA}{\mathrm{FA}}
\newcommand{\EC}{\mathrm{EC}}
\newcommand{\Vol}{\mathrm{Vol}}
\newcommand{\ind}{\mathbf{1}}
\newcommand{\cF}{\mathcal{F}}
\newcommand{\bzero}{\mathbf{0}}
\newcommand{\iotaVec}{\boldsymbol{\iota}}
\DeclareMathOperator{\E}{\mathbb{E}}
\DeclareMathOperator{\Var}{Var}
\DeclareMathOperator{\Cov}{Cov}
\DeclareMathOperator*{\plim}{plim}
\DeclareMathOperator*{\argmin}{arg\,min}

\theoremstyle{plain}
\newtheorem{proposition}{Proposition}
\newtheorem{theorem}{Theorem}
\newtheorem{lemma}{Lemma}
\theoremstyle{definition}
\newtheorem{condition}{Condition}
\newtheorem{definition}{Definition}
\theoremstyle{remark}
\newtheorem{remark}{Remark}

% Provide hyperref autoref names for theorem-like environments
\providecommand{\conditionautorefname}{Condition}
\providecommand{\propositionautorefname}{Proposition}
\providecommand{\lemmaautorefname}{Lemma}
\providecommand{\theoremautorefname}{Theorem}
\providecommand{\definitionautorefname}{Definition}
\providecommand{\remarkautorefname}{Remark}

% Allow line breaks in \texttt for long paths
\renewcommand{\ttdefault}{lmtt}
\sloppy

\begin{document}

\title{An Anderson--Rubin Test of Sectoral-Composition Optimality\\
       at the Municipality Level\\[2pt]
       \large Specification, Identification, and the Reduced-Form Equivalence}
\author{}
\date{\today}
\maketitle

\begin{abstract}
\noindent
This document is the technical reference for the Anderson--Rubin (AR) optimality
test that is the inferential object of the paper. It (i) derives the
directional-derivative condition that grounds the null
$H_0:\beta=\bzero$ in the envelope theorem; (ii) maps the shift-share
identification framework of \citet{borusyak2022quasi} (BJS) onto our political-turnover
instrument and discusses the levels-versus-changes choice formally; (iii) writes
the instruments out from primitives, including the office-specific
pre-election window and the two timing variants; (iv) shows that the AR test
of $H_0:\beta=\bzero$ is exactly the joint Wald test on the instruments in the
reduced-form GDP equation, and that this equivalence is what makes the
estimator of script~54 a weak-instrument-robust optimality test; (v) treats
the volume-channel control formally; and (vi) collects the inference questions
raised by spatial correlation in shift-share designs.
The document is intentionally compact and self-contained: all primitive
objects (firm-level owner counts, party-exposure weights, cross-office
shifts, and instruments) are defined here.
\end{abstract}

\setlength{\parskip}{6pt}
\setlength{\parindent}{12pt}

\section*{Setup and notation}\label{sec:notation}

We index municipalities by $m\in\{1,\dots,M\}$ ($M=5{,}570$), sectors by
$j\in\{1,\dots,J\}$, years by $t\in\{2002,\dots,2017\}$, offices by
$\ell\in\{\mathrm{M},\mathrm{G},\mathrm{P}\}$ (mayor, governor, president), and
parties by $p$. Lower-case Latin letters $f,g$ index firms. Let $s(m)$
denote the state of municipality $m$. Cross-office channels are indexed by
$c\in\{\mathrm{M},\,\mathrm{M{\cdot}P},\,\mathrm{M{\cdot}G},\,\mathrm{M{\cdot}G{\cdot}P}\}$.

For office $\ell$ and year $t$, let $e_\ell(t)$ denote the most recent election
year for office $\ell$ on or before $t$. The pre-election window is
\begin{equation}\label{eq:window}
  T^{\ell}_{t} \;\equiv\; \big[e_\ell(t)-4,\; e_\ell(t)-1\big] \cap [2002,2017].
\end{equation}
Within an electoral term, $T^{\ell}_t$ is constant in $t$; at cycle boundaries,
$e_\ell(t)$ advances and $T^{\ell}_{t}$ updates. The 2003 governor/president
cycle (whose pre-window is 1998--2001, fully outside the data) is dropped.

The endogenous object is the vector of \emph{sector employment shares}
\[
  s_{jmt} \;\equiv\; \frac{n_{jmt}}{n_{mt}},
  \qquad
  n_{mt} \equiv \sum_{j=1}^{J} n_{jmt},
  \qquad
  \sum_{j=1}^{J} s_{jmt} \;=\; 1,
\]
constructed from RAIS firm--municipality--year employment counts. Stacked
across $j$, $s_{mt}\equiv(s_{1mt},\dots,s_{Jmt})'\in\Delta^{J-1}$ lies on the
unit simplex; $\iotaVec' s_{mt}=1$ and $\iotaVec'\Delta s_{mt}=0$ for any change.
The outcome is municipal log real GDP, $Y_{mt}\equiv\log(\mathrm{GDP}_{mt})$, and
$\Delta Y_{mt}\equiv Y_{mt}-Y_{m,t-1}$.

We propose the following notation for BNDES sector credit shares to keep them
distinct from employment shares:
\begin{equation}\label{eq:bndes-share-notation}
  s^{\mathrm{B}}_{jmt}
  \;\equiv\;
  \frac{\mathrm{BNDES}_{jmt}}{\sum_{k}\mathrm{BNDES}_{kmt}},
\end{equation}
where $\mathrm{BNDES}_{jmt}$ denotes BNDES disbursements to sector $j$ in
municipality $m$ in year $t$. BNDES shares are the upstream mechanism through which the political
shock is transmitted to employment composition.

The total-volume control is the unit-free ratio
\begin{equation}\label{eq:vol}
  \Vol_{mt} \;\equiv\;
  \frac{\sum_{j} \mathrm{BNDES}_{jmt}}{\mathrm{GDP}_{m,t_0}},
\end{equation}
with $t_0$ the initial period of the panel ($t_0=2002$).

%% =====================================================================
\section{The optimality test: theory}\label{sec:theory}
%% =====================================================================

\paragraph{Environment.}
A municipal economy in $(m,t)$ allocates fixed factors across $J$ sectors with
the vector of shares $s_{mt} = (s_{1mt},\dots,s_{Jmt})' \in\Delta^{J-1}$.
Local output (or any local welfare proxy) is
\[
  Y_{mt} \;=\; Y\!\left(s_{mt};\,\theta_{mt}\right),
\]
where $\theta_{mt}$ collects municipality-time fundamentals
unaffected by an infinitesimal reallocation of sector shares (e.g.,
pre-existing productivity, prices, and infrastructure). We require $Y(\cdot;\theta)$ to be twice continuously
differentiable in $s$ on the relative interior of $\Delta^{J-1}$.

\paragraph{The directional-derivative condition.}
Let $\iotaVec\equiv(1,\dots,1)'\in\mathbb{R}^{J}$ denote the all-ones vector. A feasible reallocation direction is any $v\in\mathbb{R}^J$ with
$\iotaVec'v=0$, so the perturbation $s_{mt}+\varepsilon v$ remains on the
simplex for small $\varepsilon$. At an interior optimum $s^{*}_{mt}\in
\mathrm{int}(\Delta^{J-1})$\footnote{Here $\mathrm{int}(\Delta^{J-1})$ denotes
the relative interior of the simplex: all sector shares are strictly positive.
Cells with zero employment in one or more sectors are boundary points, where
the analogous optimality condition is a Kuhn--Tucker inequality rather than
the equality in~\eqref{eq:dirderiv}.} of $Y(\cdot;\theta_{mt})$, the first-order
condition is
\begin{equation}\label{eq:dirderiv}
  \frac{\mathrm{d}}{\mathrm{d}\varepsilon} Y\!\left(s^{*}_{mt}+\varepsilon v\right)
  \bigg|_{\varepsilon=0}
  \;=\;
  \sum_{j=1}^{J}
  \frac{\partial Y(s^*_{mt})}{\partial s_j} v_j
  \;=\;
  \nabla_{s} Y(s^{*}_{mt})\cdot v
  \;=\;
  0
  \qquad\text{for every } v\perp \iotaVec.
\end{equation}
Equivalently, the gradient is parallel to $\iotaVec$, the vector normal
(orthogonal) to all feasible reallocation directions, $\nabla_{s}Y\propto
\iotaVec$, i.e.\ the social marginal returns to factors are equalized across
sectors. This is the
local-optimality benchmark.

\begin{remark}[Envelope intuition]
Equation~\eqref{eq:dirderiv} is the envelope condition for the planner's
problem $\max_{s\in\Delta^{J-1}} Y(s;\theta)$: at an interior maximizer,
infinitesimal feasible perturbations have zero first-order welfare effect
\citep{samuelson1947foundations}. Rejection of~\eqref{eq:dirderiv} therefore
identifies that the observed allocation is not at an interior maximum, and
the sign pattern of the gradient prescribes the reallocation direction that
would raise $Y$.
\end{remark}

\paragraph{From the gradient to a regression.}
Let $\Delta s_{mt}\equiv s_{mt}-s_{m,t-1}$, with
$\iotaVec'\Delta s_{mt}=0$. Since $Y(\cdot;\theta_{m,t-1})$ is twice
continuously differentiable, Taylor's theorem (first-order Taylor expansion)
around $s_{m,t-1}$ gives
\begin{equation}\label{eq:taylor}
  Y(s_{mt};\theta_{m,t-1}) - Y(s_{m,t-1};\theta_{m,t-1})
  \;=\;
  \nabla_{s}Y(s_{m,t-1};\theta_{m,t-1})'\,\Delta s_{mt}
  \;+\; R_{mt},
  \qquad
  R_{mt}=O(\|\Delta s_{mt}\|^2).
\end{equation}
Observed GDP changes also include changes in fundamentals between $t-1$ and
$t$. Adding and subtracting $Y(s_{mt};\theta_{m,t-1})$ gives the exact
decomposition
\[
  Y(s_{mt};\theta_{mt}) - Y(s_{m,t-1};\theta_{m,t-1})
  =
  \Delta Y_{mt}
  =
  \underbrace{Y(s_{mt};\theta_{m,t-1})-Y(s_{m,t-1};\theta_{m,t-1})}
    _{\text{share-composition change at fixed fundamentals}}
  +
  \underbrace{Y(s_{mt};\theta_{mt})-Y(s_{mt};\theta_{m,t-1})}
    _{\text{fundamental change at the new shares}}.
\]
Substituting the Taylor expansion for the first bracket yields
\begin{equation}\label{eq:observed-change}
  \Delta Y_{mt}
  =
  \nabla_{s}Y(s_{m,t-1};\theta_{m,t-1})'\Delta s_{mt}
  +R_{mt}+D_{mt},
  \qquad
  D_{mt}\equiv
  Y(s_{mt};\theta_{mt})-Y(s_{mt};\theta_{m,t-1}).
\end{equation}

Under local optimality, the first-order composition term is zero, but observed
GDP changes need not be zero: they may still reflect changes in fundamentals
$D_{mt}$ and the approximation error $R_{mt}$.

Thus the empirical test is not whether $\Delta Y_{mt}$ is zero. It is whether the component of $\Delta s_{mt}$ induced by the instruments predicts $\Delta Y_{mt}$ after conditioning on controls. This is valid only if the instruments are orthogonal to the residual components containing $D_{mt}$, $R_{mt}$, gradient heterogeneity, and measurement error.

Equation~\eqref{eq:observed-change} and the logic described above directly motivate a first-difference equation,
\begin{equation}\label{eq:diff-struct}
  \Delta Y_{mt}
  \;=\;
  \beta'\Delta s_{mt}
  + \rho'\Delta X_{mt}
  + e^{\Delta}_{mt},
\end{equation}
where $\beta$ is interpreted as a common projected composition gradient in the
within-simplex subspace. The residual $e^{\Delta}_{mt}$ contains unabsorbed
fundamental changes $D_{mt}$, the Taylor remainder $R_{mt}$, heterogeneity in
the local gradient, and measurement error. The vector $X_{mt}$ collects the relevant controls.

\paragraph{The projected-gradient interpretation in levels.}
The directional-derivative condition~\eqref{eq:dirderiv} is a pointwise
statement: at an interior optimum, the gradient is parallel to $\iotaVec$ at
that point. The decomposition~\eqref{eq:observed-change} bridges from the
pointwise condition to a regression in $\Delta Y$ via Taylor expansion at the
cell-specific $s_{m,t-1}$. A second bridge, expanding around a common
reference allocation $\bar s$ in the relative interior of the simplex,
delivers a regression in levels:
\begin{equation}\label{eq:taylor-lvl}
  Y(s_{mt};\theta_{mt})
  \;=\;
  Y(\bar s;\theta_{mt})
  \;+\;
  \nabla_{s} Y(\bar s;\theta_{mt})'\,(s_{mt}-\bar s)
  \;+\; \tilde R_{mt},
  \qquad
  \tilde R_{mt}=O(\|s_{mt}-\bar s\|^{2}).
\end{equation}
Project the cell-level gradient $\nabla_{s} Y(\bar s;\theta_{mt})$ onto a
common within-simplex gradient $\beta$ and rearrange:
\begin{equation}\label{eq:lvl-decomp}
\begin{split}
  Y_{mt}
  \;=\;
  & \beta'\,s_{mt}
  +\underbrace{[Y(\bar s;\theta_{mt})-\beta'\bar s]}_{\text{(A) cell-specific level}} \\
  & +\underbrace{(\nabla_{s} Y(\bar s;\theta_{mt})-\beta)'(s_{mt}-\bar s)}_{\text{(B) gradient heterogeneity}}
  +\underbrace{\tilde R_{mt}}_{\text{(C) Taylor remainder}}.
\end{split}
\end{equation}
Absorbing (A)--(C) and measurement error into a single residual $e^{L}_{mt}$
yields the levels analogue of~\eqref{eq:diff-struct},
\begin{equation}\label{eq:lvl-struct}
  Y_{mt}
  \;=\;
  \beta'\,s_{mt}
  \;+\; \rho'\,X_{mt}
  \;+\; e^{L}_{mt},
\end{equation}
where $\beta$ is the same within-simplex projected composition gradient as
in~\eqref{eq:diff-struct} and $X_{mt}$ collects the relevant controls.

Equations~\eqref{eq:diff-struct} and~\eqref{eq:lvl-struct} test the same null using different residual structures. None is a-priori
small in the raw cross-section. Both nonetheless deliver the same null,
\begin{equation}\label{eq:H0}
  H_{0}:\quad \beta = \bzero,
\end{equation}
the empirical counterpart of the directional-derivative
condition~\eqref{eq:dirderiv}: at an interior local optimum, the cell-level
gradient lies in the span of $\iotaVec$, its projection onto the
within-simplex subspace $\{v:\iotaVec' v=0\}$ is zero, and with the simplex
normalization below the corresponding $\beta$ vanishes.

\begin{remark}[Why the Taylor scale is smaller than it looks]\label{rem:late}
\eqref{eq:lvl-decomp} appears to require the linearization to hold over the
full cross-sectional spread of Brazilian sectoral mixes, where
$\|s_{mt}-\bar s\|^{2}$ is large. It does not.

Municipality fixed effects absorb each muni's own average composition.
By the time the instrument acts, the cross-Brazilian dispersion is gone.
The shift-share instrument then moves $s_{mt}$ by only a few percentage
points. The linear approximation only needs to hold locally---over the
small neighborhood the instrument explores within each municipality. This
is the same demand the differences form makes.

Two interpretive consequences. \emph{(i)} When cell-level gradients are
heterogeneous, $\hat\beta$ is an Imbens--Angrist-style weighted average of
$\nabla_{s} Y(\bar s;\theta_{mt})$, weighted by IV-induced variation in
each cell. $H_{0}:\beta=\bzero$ therefore tests ``the IV-weighted average
gradient is zero,'' not a literal pointwise optimum at every $(m,t)$.
\emph{(ii)} The differences and levels specifications use different weights
and different evaluation points. Agreement between them is robustness;
disagreement locates gradient heterogeneity across the simplex.
\end{remark}

\paragraph{Simplex constraint.}
Because $\iotaVec' s_{mt}=1$ holds in every $(m,t)$ cell, the regressor matrix
of $s_{mt}$ on a constant has reduced rank. Two operationally equivalent
solutions:
\begin{enumerate}[leftmargin=*,nosep]
  \item Drop one sector $j_0$
    and interpret each $\beta_j$ as the effect of substituting one unit of
    employment share from $j_0$ to $j$.
  \item Residualize $s_{mt}$ on $\iotaVec$ (within $(m,t)$ this is degenerate)
    or, equivalently, project onto the within-simplex subspace.
\end{enumerate}
We adopt option~(1) throughout. All elements of the estimated $\hat\beta$
should be read as differences relative to $j_0$, and the AR test on the
remaining $J-1$ coefficients is invariant to the choice of $j_0$.


%% =====================================================================
\section{Identification strategy}\label{sec:identification}
%% =====================================================================

We operationalize the levels equation~\eqref{eq:lvl-struct} as a panel
projection that absorbs the cell-specific level $Y(\bar s;\theta_{mt})$
(channel~(A) of~\eqref{eq:lvl-decomp}) into municipality and year fixed
effects---valid to the extent that this object is approximately additive
in $(m,t)$ as $\alpha_m+\delta_t$---and partials out the level (volume)
channel of the political shock with the volume control $\Vol_{mt}$ defined
in~\eqref{eq:vol}. The resulting levels equation is
\begin{equation}\label{eq:struct}
  Y_{mt} \;=\; \alpha_m + \delta_t + \beta'\,s_{mt}
            + \lambda\,\Vol_{mt} + \varepsilon_{mt},
\end{equation}
where $\beta$ is the same within-simplex projected composition gradient as
in~\eqref{eq:diff-struct} and~\eqref{eq:lvl-struct}, and $\varepsilon_{mt}$
is the residual $e^{L}_{mt}$ after the FE-and-volume partialling. The volume
term enters because the same political-turnover shock that perturbs the
sectoral mix also perturbs the aggregate level of credit; $\beta$ identifies
the composition channel only after the volume channel has been partialed out.
Identification of $\beta$ further requires the instrument to be orthogonal
to the residuals---a condition formalized
through the BJS framework below. Section~\ref{sec:vol} develops the volume
control; \autoref{sec:fe} treats the remaining controls.

\begin{remark}[Why a level regression rather than a difference]
We adopt~\eqref{eq:struct} as the primary specification because each tier's alignment $\Align^{\ell}_{mpt}$ is generally constant within an electoral cycle for most municipalities and turns over at the inauguration
year (mid-cycle changes through impeachment-style processes are too rare
to drive identification). The first-difference instrument
$\Delta\Align^{\ell}_{mpt}$ is therefore a one-period pulse at
$t=e_\ell+1$ and is mechanically zero in years $\tau\in\{2,3,4\}$ of the
cycle. This risks missing the bulk of the response as the shock can produce effect any time within the term, not just at inauguration.\footnote{Another reason to prefer the levels specification is that the policy timing (BNDES allocation, firm hiring, and the resulting composition response) can unfold over the full term, not at inauguration alone.} The levels instrument
$\Align^{\ell}_{mpt}$ keeps the shock active for the entire cycle, so the
shift-share can predict $s_{mt}$ in any of the four post-election years in
which the response materializes; the differences instrument throws away,
by construction, the variation in years $\tau\in\{2,3,4\}$.
\end{remark}

\subsection{Cross-office channel structure}\label{subsec:cross-office}

The political-economy mechanism through which BNDES allocation responds
to party alignment is \emph{cross-office}---it bites when the same party (or its electoral coalition) holds power simultaneously at the mayoral office and at one or more higher tiers. The
mayor controls the local interface with BNDES (procurement, BNDES Card
take-up at municipal counters, signaling to firms about which projects
will be supported), but the bulk of BNDES disbursement decisions are made
at the regional or national level, where governors and the federal
executive have leverage. A firm in municipality $m$ aligned with party
$p$ realizes an effective political-credit advantage when both ends of
the intermediation chain align:
\begin{itemize}[leftmargin=*,nosep]
  \item the mayor of $m$ belongs to $p$ (local gatekeeping); \emph{and}
  \item either (i) the president belongs to $p$ or to a formal coalition
    partner of $p$ (national allocation favoring coalition-aligned
    municipalities), \emph{or} (ii) the governor of $s(m)$ belongs to
    $p$ or its coalition (state-level discretion over BNDES regional
    offices and over the state apparatus that complements federal
    credit programs).
\end{itemize}
A national-only or governor-only alignment that does not pass through a
sympathetic mayor lacks the local intermediary required for
sector-specific credit to land in $m$ and is not expected to drive
$s_{mt}$. We therefore restrict the instrument set to channels that
\emph{contain mayoral alignment as a factor}.

Operationally, define the cross-office shifts
\begin{align}
  \Align^{\mathrm{M{\cdot}P}}_{mpt}
    &\equiv \Align^{\mathrm{M}}_{mpt}\cdot\Align^{\mathrm{P}}_{pt},
    \label{eq:align-MP}\\
  \Align^{\mathrm{M{\cdot}G}}_{mpt}
    &\equiv \Align^{\mathrm{M}}_{mpt}\cdot\Align^{\mathrm{G}}_{s(m),pt},
    \label{eq:align-MG}\\
  \Align^{\mathrm{M{\cdot}G{\cdot}P}}_{mpt}
    &\equiv \Align^{\mathrm{M}}_{mpt}\cdot\Align^{\mathrm{G}}_{s(m),pt}
            \cdot\Align^{\mathrm{P}}_{pt},
    \label{eq:align-MGP}
\end{align}
where $\Align^{\mathrm{P}}_{pt}=1$ if party $p$ is in the federal
presidential coalition (leading party or formal coalition partner) at
time $t$, and $\Align^{\mathrm{G}}_{s(m),pt}=1$ if $p$ is in the
gubernatorial coalition of state $s(m)$ at $t$. The coalition definition
on the higher tiers reflects that BNDES allocation under any given
mandate is shared across coalition members, not monopolized by the
leading party. The pure $\Align^{\mathrm{M}}_{mpt}$ represents the standalone mayoral channel (capturing local-only
political capital, e.g.\ municipal procurement and local-bank
signaling).

\paragraph{Office-specificity of Brazilian coalitions.}
Brazilian electoral coalitions (\emph{coligações}) are formed
\emph{per office per election}: the coalition supporting a presidential
candidate is a distinct legal object from the coalition supporting a
gubernatorial candidate in any given state, and from the coalition
supporting a mayoral candidate in any given municipality. Constitutional
Amendment 52/2006 (modifying Art.\ 17, \S1) explicitly guarantees party
autonomy ``without obligatory linkage between candidacies at national,
state, district or municipal level.'' Two parties allied for a mayoral
race in muni $m$ may run on rival slates for the gubernatorial race in
state $s(m)$ in the same election cycle, and either configuration may
be unrelated to their national-presidential alignment. We therefore
treat $\Align^{\mathrm{M}}$, $\Align^{\mathrm{G}}$, and
$\Align^{\mathrm{P}}$ as separately measured objects, each defined from
the office-specific \emph{electoral coalition} (the registered
\emph{coligação} on the candidate slate of office $\ell$ at election
$e_\ell$).

Before EC 52/2006, TSE Resolution 20.993/2002 imposed
\emph{verticalização} on the 2002 and 2006 elections, requiring
state-level coalitions to align with the national presidential
coalition. Within the 2002--2017 sample window, this constraint binds
the 2002 and 2006 electoral cycles (gubernatorial inaugurations 2003 and
2007; post-election years 2003--2010).

\subsection{Firm-level primitives}\label{subsec:firm-primitives}

Here we define the firm-level primitives from which the sector-level exposures $w^{\ell}_{jmp,t}$ and firm-level instruments $\FA^{c}_{fmt}$ are built. The firm-level construction is the basis for the

\paragraph{Owner counts and firm support.}
Let $L_{f,p,t}$ denote the number of owners of firm $f$ affiliated with
party $p$ at time $t$, and $L_{f,t}$ the total number of owners of firm
$f$ at time $t$ (including owners not affiliated with any party), so
that $\sum_{p} L_{f,p,t}\le L_{f,t}$. Let $\cF(j,m)$ denote the firm
support of cell $(j,m)$: the set of firms that ever appear in sector
$j$, municipality $m$ over the panel. The pre-period firm support
associated with mayoral inauguration year $e_{\mathrm{M}}$ is
\[
  \cF^{\mathrm{pre}}_{jm,e_{\mathrm{M}}}
  \;\equiv\;
  \big\{\, f\in\cF(j,m) : L_{f,s}>0
       \text{ for some } s\in T^{\mathrm{M}}_{e_{\mathrm{M}}} \,\big\},
\]
i.e.\ the firms in $(j,m)$ active at least one year of the mayoral
pre-election window.

\paragraph{Firm-level baseline party exposure.}
The baseline exposure of firm $f$ to party $p$ for office $\ell$ at year
$t$ pools affiliated and total owner counts across the pre-election
window:
\begin{equation}\label{eq:omega-firm}
  \omega^{\ell}_{fp,t}
  \;\equiv\;
  \frac{\sum_{s\in T^{\ell}_{t}} L_{f,p,s}}
       {\sum_{s\in T^{\ell}_{t}} L_{f,s}},
  \qquad
  \sum_{p} \omega^{\ell}_{fp,t} \;\le\; 1.
\end{equation}
This pooled-count construction weights each owner-year observation
equally rather than each calendar year equally: a year in which firm
$f$ has many owners contributes more to the baseline than a year in
which it has few. Equivalently, $\omega^{\ell}_{fp,t}$ is the
owner-year-weighted average of the annual affiliation shares
$L_{f,p,s}/L_{f,s}$. The residual $1-\sum_{p}\omega^{\ell}_{fp,t}$ is
the unaffiliated-owner share at the firm level. The exposure is
office-specific because the pre-election window $T^{\ell}_{t}$ is, but
does not vary across municipalities.

\paragraph{Firm-level cross-office instrument.}
For channel $c\in\{\mathrm{M},\mathrm{M{\cdot}P},\mathrm{M{\cdot}G},
\mathrm{M{\cdot}G{\cdot}P}\}$, the firm-level cross-office instrument
is
\begin{equation}\label{eq:FA-firm}
  \FA^{c}_{fmt}
  \;\equiv\;
  \sum_{p}\,\omega^{\mathrm{M}}_{fp,t}\,\Align^{c}_{mpt},
\end{equation}
i.e.\ firm $f$'s mayoral baseline exposure inner-producted with the
channel-$c$ shift in municipality $m$ at time $t$. We use the mayoral
exposure $\omega^{\mathrm{M}}_{fp,t}$ uniformly across channels; the
substantive justification is in \autoref{rem:weights} below.

\paragraph{Sector-level baseline exposure.}
The sector-level baseline exposure aggregates firm-level
exposures~\eqref{eq:omega-firm} across firms within $(j,m)$ using
owner-count weights, with the office-$\ell$ pre-election window:
\begin{equation}\label{eq:wjmp}
  w^{\ell}_{jmp,t}
  \;\equiv\;
  \frac{1}{|T^{\ell}_t|}\sum_{s\in T^{\ell}_t}
    \frac{\sum_{f\in\cF(j,m)} L_{f,p,s}}
         {\sum_{f\in\cF(j,m)} L_{f,s}},
  \qquad
  \sum_{p} w^{\ell}_{jmp,t}\le 1.
\end{equation}
The residual $1-\sum_{p}w^{\ell}_{jmp,t}$ is the unaffiliated-owner
share of cell $(j,m)$, which contributes zero to $Z^{\ell}_{jmt}$ but
enters the exposure control
$\EC^{\ell}_{jm,t}\equiv\sum_p w^{\ell}_{jmp,t}$ that absorbs total
political connectedness of $(j,m)$.

\paragraph{Owner-count normalizer.}
Let
\begin{equation}\label{eq:N-norm}
  N^{\ell}_{jm,t}
  \;\equiv\;
  \sum_{f\in\cF(j,m)} L_{f,t},
\end{equation}
the total number of owners across firms in cell $(j,m)$ at time $t$.

\paragraph{Sector exposure as owner-count-weighted firm average.}
Substituting~\eqref{eq:omega-firm} into~\eqref{eq:wjmp} and rearranging
gives the equivalent representation
\begin{equation}\label{eq:wjmp-as-avg}
  w^{\ell}_{jmp,t}
  \;=\;
  \sum_{f\in\cF(j,m)}
       \frac{L_{f,t}}{N^{\ell}_{jm,t}}\,\omega^{\ell}_{fp,t}.
\end{equation}
The sector-level baseline party exposure is the owner-count-weighted
average of firm-level exposures across the firm support
$\cF(j,m)$.\footnote{The simple (equal-firm-weight) average
$\bar\omega^{\ell}_{jmp,t}\equiv |\cF(j,m)|^{-1}\sum_{f}\omega^{\ell}_{fp,t}$
differs from~\eqref{eq:wjmp-as-avg} whenever firm-level exposure
$\omega^{\ell}_{fp,t}$ covaries with firm size $L_{f,t}$. Owner-count
weighting matches the empirical mass of political exposure in the cell
and is the construction used throughout the pipeline.}
Equation~\eqref{eq:wjmp-as-avg} is the bridge between the firm-level
instrument~\eqref{eq:FA-firm} and the sector-level instruments of the
next subsection.

\subsection{The shift-share form of the cross-office instruments}

Building on the firm-level instrument~\eqref{eq:FA-firm} and the
sector-level baseline exposure~\eqref{eq:wjmp}, the four sector-level
levels instruments are
\begin{align}
  Z^{\mathrm{M}}_{jmt}
    &\equiv \sum_{p} w^{\mathrm{M}}_{jmp,t}\,\Align^{\mathrm{M}}_{mpt},
    \label{eq:Z-M}\\
  Z^{\mathrm{M{\cdot}P}}_{jmt}
    &\equiv \sum_{p} w^{\mathrm{M}}_{jmp,t}\,\Align^{\mathrm{M{\cdot}P}}_{mpt},
    \label{eq:Z-MP}\\
  Z^{\mathrm{M{\cdot}G}}_{jmt}
    &\equiv \sum_{p} w^{\mathrm{M}}_{jmp,t}\,\Align^{\mathrm{M{\cdot}G}}_{mpt},
    \label{eq:Z-MG}\\
  Z^{\mathrm{M{\cdot}G{\cdot}P}}_{jmt}
    &\equiv \sum_{p} w^{\mathrm{M}}_{jmp,t}\,\Align^{\mathrm{M{\cdot}G{\cdot}P}}_{mpt}.
    \label{eq:Z-MGP}
\end{align}
For backward compatibility with the rest of the document we let
$Z^{c}_{jmt}$ denote a generic channel
$c\in\{\mathrm{M},\mathrm{M{\cdot}P},\mathrm{M{\cdot}G},
\mathrm{M{\cdot}G{\cdot}P}\}$ and $\Align^{c}_{mpt}$ its shift. The
changes variant is
\begin{equation}\label{eq:dZ}
  \Delta Z^{c}_{jmt}
  \;\equiv\;
  \sum_{p} w^{\mathrm{M}}_{jmp,t}\,\Delta\Align^{c}_{mpt},
  \qquad
  \Delta\Align^{c}_{mpt} \equiv \Align^{c}_{mpt}-\Align^{c}_{mp,t-1},
\end{equation}
which is non-zero in any year in which any of the offices entering $c$
turns over.

\begin{remark}[Mayoral exposure for all channels]\label{rem:weights}
We use the mayoral exposure $w^{\mathrm{M}}_{jmp,t}$ for all four
cross-office instruments. The mayor is the local intermediary in every
channel, so the sector-municipality exposure to ``politically active
firms'' is naturally a mayor-relative object. A coalition-specific
alternative $w^{c}_{jmp,t}$ that counts only owners affiliated with
parties in the relevant federal/state coalition during the mayoral
pre-election window would concentrate exposure on the channel-relevant
subset and possibly improve first-stage power. The cost is a more
complex \autoref{cond:BJS1}: the exposure shares would inherit a
presidential- or gubernatorial-cycle dependence that the mayor-only
$w^{\mathrm{M}}$ does not have. We adopt $w^{\mathrm{M}}$ as the primary
specification for cleanliness of identification and treat the
coalition-specific construction as a robustness exercise.
\end{remark}

\paragraph{The owner-count-weighted aggregation identity.}
Each $Z^{c}_{jmt}$ is a within-region inner product of \emph{shares}
$w^{\mathrm{M}}_{jmp,t}$ (sector-municipality exposure to party $p$, on
the mayoral pre-election window) and \emph{shifts} $\Align^{c}_{mpt}$
(whether $p$ controls the offices defining $c$ at time $t$). It is the
sector-aggregated counterpart of the firm-level
instrument~\eqref{eq:FA-firm}: substituting~\eqref{eq:wjmp-as-avg}
into~\eqref{eq:Z-M}--\eqref{eq:Z-MGP} and reordering sums gives
\begin{equation}\label{eq:agg-identity}
  Z^{c}_{jmt}
  \;=\;
  \sum_{p} w^{\mathrm{M}}_{jmp,t}\,\Align^{c}_{mpt}
  \;=\;
  \sum_{p}\Bigg[\sum_{f\in\cF^{\mathrm{pre}}_{jm,e_{\mathrm{M}}}}
       \frac{L_{f,t}}{N^{\mathrm{M}}_{jm,t}}\,\omega^{\mathrm{M}}_{fp,t}\Bigg]
       \Align^{c}_{mpt}
  \;=\;
  \sum_{f\in\cF^{\mathrm{pre}}_{jm,e_{\mathrm{M}}}}
       \frac{L_{f,t}}{N^{\mathrm{M}}_{jm,t}}\,\FA^{c}_{fmt}.
\end{equation}
The sector instrument is a weighted average of the firm-level
instruments across firms in the pre-period support, with weights given
by each firm's share of total owner-counts in the cell. The
identity~\eqref{eq:agg-identity} verifies diagnostically: collapsing
firm instruments by these weights recovers the sector instrument.

\subsection{Mass-weighted exposure alternatives}\label{subsec:mass-weighted}

The intensity-based exposure $w^{\mathrm{M}}_{jmp,t}$ in~\eqref{eq:wjmp}
normalizes owner-party counts \emph{within} each $(j,m)$ cell, so it
captures the intensity of political affiliation in the cell but discards
information about the cell's mass in the local economy. For predicting
sector employment shares $s_{jmt}$ this is a known weakness of the
intensity-only shift-share~\eqref{eq:Z-M}--\eqref{eq:Z-MGP}: a small cell
with a high aligned-fraction generates a large instrument value but can
move little employment when credit arrives, while a large cell with a
moderate aligned-fraction can move a lot. This subsection records two
alternative exposure specifications, both predetermined on
$T^{\mathrm{M}}_{t}$ and both preserving the shift-share decomposition,
that bring mass back into the share component.

\paragraph{Pre-window employment primitives.}
Let $E_{f,t}$ denote firm $f$'s employment at time $t$. Define the
pre-window firm employment, sector-muni pre-window employment, and muni
pre-window employment as the natural averages over the mayoral
pre-election window $T^{\mathrm{M}}_t$:
\begin{equation}\label{eq:Ebar}
  \bar E_{f,t}\;\equiv\;
   \frac{1}{|T^{\mathrm{M}}_t|}\sum_{s\in T^{\mathrm{M}}_t} E_{f,s},
  \qquad
  \bar E_{jm,t}\;\equiv\;\sum_{f\in\cF(j,m)} \bar E_{f,t},
  \qquad
  \bar E_{m,t}\;\equiv\;\sum_{j'} \bar E_{j'm,t}.
\end{equation}
The pre-window matches the one used for owner-party counts
in~\eqref{eq:omega-firm} and~\eqref{eq:wjmp}, so the predetermination
argument that justifies $w^{\mathrm{M}}_{jmp,t}$ extends without
modification to any object built from $\bar E$.

\paragraph{Variant A: muni-relative aligned-owner share.}
The first variant keeps the owner-count primitives but replaces the
within-cell normalization in~\eqref{eq:wjmp} with a muni-level
normalization over affiliated owners. Let
\(
  \bar L_{jmp,t}\equiv\sum_{s\in T^{\mathrm{M}}_t}
    \sum_{f\in\cF(j,m)} L_{f,p,s}
\)
be the pre-window owner-years in $(j,m)$ affiliated with party $p$, and
\(
  \bar L^{\mathrm{affil}}_{m,t}\equiv\sum_{j'}\sum_{p'}\bar L_{j'mp',t}
\)
the muni-level pre-window total of owner-years affiliated with any party.
Define
\begin{equation}\label{eq:w-own-rel}
  \tilde w^{\mathrm{own}}_{jmp,t}
   \;\equiv\;
   \frac{\bar L_{jmp,t}}{\bar L^{\mathrm{affil}}_{m,t}},
   \qquad
   \tilde Z^{c,\mathrm{own}}_{jmt}
   \;\equiv\;
   \sum_p \tilde w^{\mathrm{own}}_{jmp,t}\,\Align^{c}_{mpt}.
\end{equation}
The share $\tilde w^{\mathrm{own}}_{jmp,t}$ is the fraction of the
muni's politically-engaged pre-window owner-mass that sits in sector
$j$ at party $p$; it carries information about cell size through
$\bar L_{jmp,t}$ that the intensity $w^{\mathrm{M}}_{jmp,t}$ discards.
Both numerator and denominator are pre-window quantities, so
$\tilde w^{\mathrm{own}}_{jmp,t}$ is channel-independent (consistent
with~\autoref{rem:weights}) and predetermined relative to the shift
$\Align^{c}_{mpt}$. Across sectors, $\sum_j \tilde Z^{c,\mathrm{own}}_{jmt}
\in[0,1]$ equals the muni's share of pre-window affiliated owners that
lean toward coalition-$c$ parties.

\paragraph{Variant B: muni-relative aligned-employment share.}
The second variant replaces owner-mass by employment-mass at every
aggregation level. Using the firm-level
instrument~\eqref{eq:FA-firm}, define
\begin{equation}\label{eq:Z-emp-rel}
  \tilde Z^{c,\mathrm{emp}}_{jmt}
   \;\equiv\;
   \frac{\displaystyle\sum_{f\in\cF(j,m)} \bar E_{f,t}\,\FA^{c}_{fmt}}
        {\bar E_{m,t}}.
\end{equation}
This is the pre-window employment of firms in $(j,m)$ whose owners lean
toward coalition-$c$ parties (with $\FA^{c}_{fmt}\in[0,1]$ as the
per-firm alignment-intensity weight), expressed as a fraction of the
muni's total pre-window employment. The same object admits the
shift-share decomposition
\begin{equation}\label{eq:Z-emp-rel-ss}
  \tilde Z^{c,\mathrm{emp}}_{jmt}
   \;=\;
   \sum_p \tilde w^{\mathrm{emp}}_{jmp,t}\,\Align^{c}_{mpt},
   \qquad
   \tilde w^{\mathrm{emp}}_{jmp,t}
   \;\equiv\;
   \frac{\displaystyle\sum_{f\in\cF(j,m)} \bar E_{f,t}\,\omega^{\mathrm{M}}_{fp,t}}
        {\bar E_{m,t}}.
\end{equation}
The share admits a size-times-intensity decomposition
\begin{equation}\label{eq:w-emp-decomp}
  \tilde w^{\mathrm{emp}}_{jmp,t}
   \;=\;
   \underbrace{\frac{\bar E_{jm,t}}{\bar E_{m,t}}}_{\text{sector size}}\;\;
   \underbrace{\bar\omega^{\mathrm{M},E}_{jmp,t}}_{\text{employment-weighted intensity}},
   \qquad
   \bar\omega^{\mathrm{M},E}_{jmp,t}\;\equiv\;
   \sum_{f\in\cF(j,m)} \frac{\bar E_{f,t}}{\bar E_{jm,t}}\,
   \omega^{\mathrm{M}}_{fp,t},
\end{equation}
in which $\bar\omega^{\mathrm{M},E}_{jmp,t}$ is the
employment-weighted within-cell average of firm-level intensities ---
the direct analogue of~\eqref{eq:wjmp-as-avg} with owner-count weights
replaced by employment weights. Both factors are pre-window and
channel-independent; the shift sits exclusively in
$\Align^{c}_{mpt}$ as in the base spec.

\paragraph{Comparison of A and B.}
The variants address different facets of the mass-weighting concern.
Variant~A modifies the across-sector scaling: muni-level political
mass is allocated to sectors in proportion to where the aligned
\emph{owners} are. Variant~B modifies \emph{both} the firm-to-sector
aggregation~\eqref{eq:wjmp-as-avg} and the across-sector scaling,
replacing owner-mass by \emph{employment-mass} everywhere. The two
coincide in within-muni ranking only when firm-level employment is
proportional to owner counts; in industry-towns, where a small number
of large employers carry most of the employment but a minority of the
owners, Variant~B will rank sectors differently from Variant~A and
from the base intensity spec. Variant~B is preferred on the
mechanism-fit argument that the conversion of credit to hires
operates through firm size in employment, not in owners; Variant~A is
the conservative analogue that preserves the owner-count primitives
already used throughout~\eqref{eq:omega-firm}--\eqref{eq:wjmp-as-avg}.

\paragraph{BJS conditions.}
\autoref{cond:BJS1} and~\autoref{cond:BJS2} carry through for both
variants under the same argument as~\eqref{eq:Z-M}--\eqref{eq:Z-MGP}:
the shift $\Align^{c}_{mpt}$ is unchanged, and the share components
are pre-window objects.~\autoref{cond:BJS3} is the condition where
the variants can deviate from the base spec, because muni-relative
normalization concentrates exposure on large-mass cells; the
Herfindahl $\sum_j (\tilde Z^{c,\bullet}_{jmt})^2$ should be reported
alongside the base-spec Herfindahl in robustness diagnostics, since
industry-towns load disproportionately on a single sector under
Variants~A and B.

\paragraph{Exposure controls.}
The base-spec exposure control $\EC^{\ell}_{jm,t}\equiv\sum_p
w^{\ell}_{jmp,t}$ adapts mechanically:
$\widetilde\EC^{\mathrm{own}}_{jm,t}\equiv
\sum_p \tilde w^{\mathrm{own}}_{jmp,t}$ is cell $(j,m)$'s share of
muni affiliated owners under Variant~A, and
$\widetilde\EC^{\mathrm{emp}}_{jm,t}\equiv
\sum_p \tilde w^{\mathrm{emp}}_{jmp,t}$ its employment-mass analogue
under Variant~B. The muni-level political-mass scales
$\log\bar L^{\mathrm{affil}}_{m,t}$ (Variant~A) and $\log\bar E_{m,t}$
(Variant~B) are added to the controls in~\eqref{eq:struct} to absorb
the level of muni political mass without absorbing the across-sector
tilt that identifies $\beta$.

\paragraph{Implementation.}
We treat the intensity-based
instruments~\eqref{eq:Z-M}--\eqref{eq:Z-MGP} as the base specification
and report $\tilde Z^{c,\mathrm{own}}_{jmt}$ and
$\tilde Z^{c,\mathrm{emp}}_{jmt}$ alongside in first-stage and AR
robustness tables. The three specifications are non-nested but share
the same identification skeleton; differences across them are
informative about which mass concept actually predicts sector
employment shares in the data.

\subsection{BJS conditions, mapped to our setting}

\citet{borusyak2022quasi} treat shift-share IVs as IV estimators in a
\emph{shock-level} sample, where the unit of identification is a
$(\text{shock})$ rather than an $(m,t)$ cell. In our notation the
``shocks'' are office-party-municipality cells $k\equiv(\ell,p,m,t)$ with
shifts $g_k\equiv\Align^{\ell}_{mpt}$ (or $\Delta\Align^{\ell}_{mpt}$) and
exposures $w^{\ell}_{jmp,t}$. We restate their three primitive identification
assumptions and map each to our setting.

\begin{condition}[BJS-1: Quasi-random shock assignment]\label{cond:BJS1}
Conditional on a vector of shock-level controls $q_k$,
\(
  \E[g_k \mid \{w^{\ell}_{jmp,t}\}_{j},\, q_k] = \mu(q_k)
\)
is the same for shocks with different exposure profiles. Equivalently, in our
panel: among $(\ell,p,m,t)$ cells with the same controls $q_k$, who is aligned
is as good as randomly assigned with respect to municipal GDP residuals.
\end{condition}

\begin{condition}[BJS-2: Many uncorrelated shocks]\label{cond:BJS2}
The number of effective shocks $K_{\mathrm{eff}}$ is large, and shocks are
uncorrelated with each other after conditioning on $q_k$.
\end{condition}

\begin{condition}[BJS-3: Non-concentration of average exposure]\label{cond:BJS3}
The Herfindahl of the average exposure weights across shocks is small, so no
single shock drives the IV estimate.
\end{condition}

\paragraph{What each buys.}
\autoref{cond:BJS1} is the analogue of the exclusion restriction. It is what
allows the instrument $Z^{\ell}_{jmt}$ to be exogenous in the structural
GDP equation. \autoref{cond:BJS2} provides the law-of-large-numbers justification
that shock-level orthogonality of $g_k$ to municipal residuals translates into
unit-level (muni-year) orthogonality of $Z^{\ell}_{jmt}$ to $\varepsilon_{mt}$.
\autoref{cond:BJS3} ensures that the asymptotic distribution of the IV
estimator is non-degenerate and not driven by a handful of large-exposure
shocks.

\paragraph{What can violate \autoref{cond:BJS1} in our setting.}
The political-alignment shock $\Align^{\ell}_{mpt}$ is not literally random:
parties win office because voters prefer them. Quasi-random assignment is
plausible only after conditioning on a sufficient set $q_k$ of shock-level
controls. The relevant candidates for $q_k$ are: (a) state-by-year fixed
effects (so identification is within-state contests); (b) lagged municipal
political controls (a party's pre-period vote share); (c) lagged GDP and
employment trends (so a party that systematically wins in growing
municipalities is not confounded with growth); and (d) sector-by-year fixed
effects in the first stage (\autoref{sec:firststage}) so that nationwide
sectoral cycles are absorbed at the shock level. The empirical strategy
includes muni and year fixed effects, exposure controls $\EC^{\ell}_{jm,t}$,
and (in robustness) a pre-period GDP trend; see \autoref{sec:fe}.

\subsection{Levels versus changes: a formal sufficient condition}

The same shift-share construction admits two timing variants of the
shift, applied to each channel
$c\in\{\mathrm{M},\mathrm{M{\cdot}P},\mathrm{M{\cdot}G},\mathrm{M{\cdot}G{\cdot}P}\}$:
$\Align^{c}_{mpt}\in\{0,1\}$ in levels, or
$\Delta\Align^{c}_{mpt}\in\{-1,0,+1\}$ in changes. They differ in how the
identifying variation is distributed over time. Let $E_c$ denote the set
of inauguration years for any office defining $c$:
$E_{\mathrm{M}}=\{2005,2009,2013,2017\}$, while
$E_{\mathrm{M{\cdot}P}}=E_{\mathrm{M{\cdot}G}}=
E_{\mathrm{M{\cdot}G{\cdot}P}}=\{2003,2005,2007,2009,\ldots,2017\}$ is
the union of mayoral and gubernatorial/presidential inaugurations.

\begin{proposition}[Sufficient conditions for the levels and changes
instruments]\label{prop:lvlchg}
Suppose $Y_{mt}=\alpha_m+\delta_t+\beta' s_{mt}+\lambda\Vol_{mt}+\varepsilon_{mt}$
with $\E[\varepsilon_{mt}\mid \alpha_m,\delta_t]=0$, and let
$\mathcal{H}_{m,t-1}\equiv\sigma(Y_{m,t'},\Vol_{m,t'},s_{m,t'};t'<t)$ be the
pre-period information set.

(i) The \emph{levels} instrument $Z^{c}_{jmt}$ satisfies the BJS
quasi-random assignment condition \autoref{cond:BJS1} if, between
consecutive inauguration dates $e,e'\in E_c$ with $e<e'$,
\begin{equation}\label{eq:lvl-cond}
  \E\!\left[\Align^{c}_{mp,e+\tau}\,\big|\,
    \{w^{\mathrm{M}}_{jmp,e}\}_{j},\,\mathcal{H}_{m,e-1},\, q_{k}\right]
  \;=\;
  \mu_\tau(q_k),
  \qquad \tau=1,\dots,e'-e,
\end{equation}
i.e.\ the channel-$c$ alignment status is uncorrelated with prior
municipal trajectories conditional on $q_k$, throughout the entire
post-inauguration interval.

(ii) The \emph{changes} instrument $\Delta Z^{c}_{jmt}$ satisfies
\autoref{cond:BJS1} if, for every $t\in E_c$,
\begin{equation}\label{eq:chg-cond}
  \E\!\left[\Delta\Align^{c}_{mpt}\,\big|\,
    \{w^{\mathrm{M}}_{jmp,t}\}_{j},\,\mathcal{H}_{m,t-1},\, q_{k}\right]
  \;=\;
  \mu(q_k),
\end{equation}
i.e.\ each channel-$c$ turnover is uncorrelated with prior trajectories
conditional on $q_k$.
\end{proposition}

\paragraph{Discussion.}
Equation~\eqref{eq:lvl-cond} requires conditional mean-independence of
the channel-$c$ alignment from pre-period outcomes \emph{at every horizon}
within the relevant inter-inauguration interval.
Equation~\eqref{eq:chg-cond} requires it only \emph{at the moments of
turnover}. The latter is empirically weaker: turnover is closer to a
discontinuity at each election date and need only be unrelated to the
trajectory entering that election, not to the entire path that follows.

\paragraph{Inheritance from the product structure.}
For the cross-office channels $c\in\{\mathrm{M{\cdot}P},
\mathrm{M{\cdot}G},\mathrm{M{\cdot}G{\cdot}P}\}$ the shift
$\Align^{c}_{mpt}$ is a product of tier-specific shifts. A sufficient
primitive condition for~\eqref{eq:lvl-cond}--\eqref{eq:chg-cond} is that
each component $\Align^{\mathrm{M}}_{mpt}$, $\Align^{\mathrm{G}}_{s(m),pt}$,
$\Align^{\mathrm{P}}_{pt}$ separately satisfies the corresponding
mean-independence condition \emph{and} that the components are
conditionally independent given $q_k$:
\begin{equation}\label{eq:prod-indep}
  \E\!\left[\Align^{\mathrm{M}}_{mpt}\,\Align^{\mathrm{P}}_{pt}\,\big|\,
    \mathcal{H}_{m,t-1},q_k\right]
  \;=\;
  \E\!\left[\Align^{\mathrm{M}}_{mpt}\,\big|\,
    \mathcal{H}_{m,t-1},q_k\right]\cdot
  \E\!\left[\Align^{\mathrm{P}}_{pt}\,\big|\,q_k\right],
\end{equation}
and analogously for $\Align^{\mathrm{G}}$ and the triple product.
Substantively this rules out strategic anticipation: for instance, the
mayoral election outcome in $m$ must not be systematically aligned with
the federal incumbent's coattail strength beyond what $q_k$ already
absorbs. Including state$\times$year fixed effects and pre-period vote
shares in $q_k$ is the natural way to discharge this condition.

\paragraph{Verticalização caveat for cycles 2002 and 2006.}
The product-independence assumption~\eqref{eq:prod-indep} is partially
violated for the 2002 and 2006 electoral cycles by TSE's
\emph{verticalização} rule (Resolution 20.993/2002, in force until EC
52/2006 took effect for the 2010 elections). During those cycles, state
gubernatorial coalitions were required to mirror the national presidential
coalition, so $\Align^{\mathrm{G}}_{s(m),pt}$ and $\Align^{\mathrm{P}}_{pt}$
are not conditionally independent in the BJS sense: their joint distribution
collapses toward perfect alignment by legal construction. The mechanical
consequence is that $Z^{\mathrm{M{\cdot}G}}_{jmt}$ and $Z^{\mathrm{M{\cdot}P}}_{jmt}$
are more highly correlated for cycles 2002 and 2006 than for cycles 2010
and 2014, and the joint AR test on the four-channel cross-office stack pools
verticalizado and non-verticalizado identifying variation. Two practical
responses: (i) report the joint AR statistic with a footnote flagging
the verticalização period; (ii) report a robustness AR test restricted
to cycles 2010 and 2014 (post-verticalização sample), where
\eqref{eq:prod-indep} is constitutionally protected (EC 52/2006).

\paragraph{Within-cycle persistence (levels).}
In the levels specification for the pure mayoral channel,
$\Align^{\mathrm{M}}_{mpt}$ is non-zero for the entire $\sim 4$-year
mayoral cycle, so $q_k$ must absorb anything that correlates persistently
within a cycle with both alignment and outcomes (e.g.\ a slow-moving
regional growth differential). Including a pre-period GDP trend, lagged
$Y_{m,t-1}$, or muni-specific linear trends are natural strengthenings;
the 2002-fixed baseline in $\omega^{\mathrm{M}}_{fp,t}$ also reduces the
dependence on contemporaneous alignment-cycle dynamics. For the cross-office
channels, the inter-inauguration intervals are at most two years, so the
within-interval persistence requirement is correspondingly weaker: the
levels-versus-changes gap shrinks for $c\in\{\mathrm{M{\cdot}P},
\mathrm{M{\cdot}G},\mathrm{M{\cdot}G{\cdot}P}\}$.

\paragraph{Many-shocks count under \autoref{cond:BJS2}.}
The effective number of shocks differs sharply between the two variants.
With $M=5{,}570$, four mayoral cycles
($e_{\mathrm{M}}\in\{2004,2008,2012,2016\}$) and $P\approx 30$ parties,
the total \emph{candidate} shock count is on the order of $M\times P\times
4\approx 6.7\times 10^{5}$ for the pure mayoral channel. The levels
variant uses each shock for $\sim 4$ post-election years, but for the
pure mayoral channel those four year-replicates contribute redundant
information for shock-level inference (the shift is constant within a
cycle). The changes variant uses each shock once at inauguration, so its
effective shock count is the same but its panel-level non-zero cell count
is smaller by a factor of $\sim 4$. For the cross-office channels the
``constant within cycle'' redundancy is already broken---the shift can
flip mid-mayoral-term at the gubernatorial/presidential
inauguration---so the levels-variant cell count is roughly $2\times$
rather than $4\times$ the changes-variant count, and the levels-versus-changes
shock-level redundancy is correspondingly smaller.
\autoref{cond:BJS2} is therefore satisfied at comparable levels by both
variants in terms of \emph{distinct shocks}, but the changes variant has less
redundancy and a more transparent shock-level interpretation. By contrast, the
levels variant has more non-zero panel cells and thus more degrees of freedom
for the muni-year regression--at the cost of a stronger \autoref{cond:BJS1}
requirement (\eqref{eq:lvl-cond}).

\paragraph{Operational rule.}
Levels instruments are the primary specification, on grounds of statistical
power and continuity with the existing pipeline. The changes instruments are
the robustness specification. If the two reject in agreement, the
substantive conclusion is robust to whichever of \eqref{eq:lvl-cond} or
\eqref{eq:chg-cond} actually holds in the data; if they disagree, the
discrepancy isolates exactly the part of the variation that
\eqref{eq:lvl-cond} is asking about beyond \eqref{eq:chg-cond}.

\subsection{Exclusion restriction}

The exclusion restriction proper is that the instrument
$Z^{\ell}_{jmt}$ enters $Y_{mt}$ only through the
sectoral employment composition $s_{mt}$, after conditioning on the volume
control $\Vol_{mt}$ and on the controls of \autoref{sec:fe}. Three
non-composition channels are conceivable and are addressed in turn:
\begin{enumerate}[leftmargin=*,nosep]
\item \emph{Volume of BNDES credit.} Political alignment changes both the
  composition and the level of BNDES credit. The level effect is partialled out
  by $\Vol_{mt}$ in~\eqref{eq:struct}. The exclusion restriction asks that,
  conditional on $\Vol_{mt}$, $Z^{\ell}_{jmt}$ does not drive $Y_{mt}$ through a
  residual volume channel.
\item \emph{Non-BNDES political channels.} A new mayor may also change
  procurement, transfers, or local public investment in ways that move
  GDP. \autoref{cond:BJS1} accommodates this only if those channels are
  orthogonal to the \emph{sector-specific} pattern of $w^{\ell}_{jmp,t}$. We
  test this in robustness as a placebo: if $Z^{\ell}_{jmt}$ predicts municipal
  transfers or procurement after conditioning, the exclusion restriction is
  weakened.
\item \emph{Direct partisanship effects on private behaviour.} Even absent
  any policy change, owners aligned with the new ruling party may invest more
  in their own firms because of confidence shocks. This would deliver a
  reduced-form effect on $s_{mt}$ via firm-level investment that is not
  mediated by BNDES credit, but it is still a \emph{composition} channel and
  therefore part of the estimand by construction--it is not an exclusion
  violation. The estimand is the GDP effect of compositional shifts driven by
  the political shock, whether transmitted through BNDES credit or through
  partisan-firm investment.
\end{enumerate}

%% =====================================================================
\section{Instrument construction}\label{sec:instrument}
%% =====================================================================

This section assembles the sector-channel
instruments~\eqref{eq:Z-M}--\eqref{eq:Z-MGP} into the muni-year stack
used by the AR test, and documents the cycle-alignment structure of
the exposure weights.

\paragraph{Sector instruments and muni-level stacking.}
Equations~\eqref{eq:Z-M}--\eqref{eq:dZ} define the sector-level cross-office
instruments. To take them to the muni-year level for the AR test, we
\emph{stack} sector and channel:
\begin{equation}\label{eq:stacked}
  Z_{mt}
  \;\equiv\;
  \big( Z^{c}_{jmt} \big)_{c\in\mathcal{C},\,j=1,\dots,J}
  \;\in\;\mathbb{R}^{K},
  \qquad
  K \;=\; |\mathcal{C}|\cdot J,
\end{equation}
where $\mathcal{C}\subseteq\{\mathrm{M},\mathrm{M{\cdot}P},
\mathrm{M{\cdot}G},\mathrm{M{\cdot}G{\cdot}P}\}$ is the active channel
set. For the currently leading post-D28 candidate margin
$\texttt{policy\_block\_active}\times\texttt{S3}$ with $J=12$ active bins,
$K=48$ for the full four-channel cross-office stack
($\mathcal{C}=\{\mathrm{M},\mathrm{M{\cdot}P},\mathrm{M{\cdot}G},
\mathrm{M{\cdot}G{\cdot}P}\}$), $K=36$ for the three-channel cross-office
stack omitting the triple ($\mathcal{C}=\{\mathrm{M},\mathrm{M{\cdot}P},
\mathrm{M{\cdot}G}\}$), and $K=12$ for the pure-mayoral specification
($\mathcal{C}=\{\mathrm{M}\}$). At the secondary candidate margin
$\texttt{cnae\_section}\times\texttt{S3}$ with $J=51$ active bins,
$K\in\{51,153,204\}$ for the three nested specifications.

\paragraph{Cycle alignment and exposure weights.}
Mayoral elections and gubernatorial/presidential elections sit on
disjoint four-year cycles ($e_{\mathrm{M}}\in\{2004,2008,2012,2016\}$
versus $e_{\mathrm{G}}=e_{\mathrm{P}}\in\{2002,2006,2010,2014\}$).
Because governors and the president share the same election year, the
pre-election windows $T^{\mathrm{G}}_{t}$ and $T^{\mathrm{P}}_{t}$
defined in~\eqref{eq:window} coincide and so $w^{\mathrm{G}}_{jmp,t}
=w^{\mathrm{P}}_{jmp,t}$ when both are constructed on the common
four-year window. Within any mayoral term, exactly one
gubernatorial/presidential turnover occurs (at $t=e_{\mathrm{M}}+2$), so
each cross-office shift has at most two within-mayoral-term changes: one at
the mayoral inauguration $e_{\mathrm{M}}+1$ and one at the
gubernatorial/presidential inauguration $e_{\mathrm{M}}+3$.

\begin{remark}[Channel-set selection]\label{rem:channel-set}
The four-channel stack is the most aggressive specification and is the
natural primary if the triple-cross-office shift $\Align^{\mathrm{M{\cdot}G{\cdot}P}}$
has enough non-zero panel cells to identify $\beta$. The three-channel
stack drops the triple in case it is too thin (full-coalition alignment
is the rarest configuration: it requires the firm-affiliated party to
hold mayor, governor, and the federal coalition simultaneously). The
pure-mayoral specification is the conservative comparison and continuity
with the existing pipeline.
\end{remark}

\begin{remark}[Wide vs.\ aggregated muni-level instrument]
An alternative aggregation, ``stack-and-summarize,'' would collapse
$Z^{c}_{jmt}$ to a single muni-year object via
$Z^{c}_{mt}=\sum_j w_j Z^{c}_{jmt}$ for some weights $w_j$. This reduces
$K$ but imposes an aggregation that pre-loads which sectors matter; we
treat it as a robustness check. The wide stacking~\eqref{eq:stacked}
keeps the cross-sector pattern of the shock visible to the test and
makes the ``$\beta=\bzero$'' null intelligible coefficient by coefficient.
\end{remark}

%% =====================================================================
\section{First stage}\label{sec:firststage}
%% =====================================================================

The first stage relating sectoral employment shares to the instruments is
estimated on the muni-sector-year panel (Panel A) at the chosen aggregation
margin:
\begin{equation}\label{eq:fs}
  s_{jmt}
  \;=\;
  \sum_{\ell}\,\pi^{\ell}_j\, Z^{\ell}_{jmt}
  \;+\;
  \sum_{\ell}\,\phi^{\ell}_j\,\EC^{\ell}_{jm,t}
  \;+\;
  \gamma_{jm} + \alpha_{jt}
  \;+\;
  \eta_{jmt},
\end{equation}
with municipality$\times$sector fixed effects $\gamma_{jm}$
($\sim$\texttt{muni\_id\^{}cnae\_section} in fixest), sector$\times$year fixed
effects $\alpha_{jt}$ ($\sim$\texttt{cnae\_section\^{}year}), and standard
errors two-way clustered by $\texttt{muni\_id}$ and sector. The exposure
control $\EC^{\ell}_{jm,t}$ absorbs the total pre-period political
connectedness of $(j,m)$ across all parties and is the level-restricted
analogue of \citeauthor{goldsmith2020bartik}'s share-level controls
\citep{goldsmith2020bartik}.

The first-stage F-statistic, computed via
$\mathtt{fixest::wald(\cdot,\,keep="\char`\^Z\_")}$ on the full set of
sector-tier instruments, is the diagnostic for $\autoref{cond:BJS3}$ and for
F2 informativeness in our identification chain. Per D20 it is not a validity
gate for the AR test: weak first stages widen the AR confidence set but do
not bias inference.

\begin{remark}[Stacked first stage for the AR test]
The first stage~\eqref{eq:fs} need not be estimated separately to conduct the
AR test: as the next section shows, it is internally embedded in the
reduced-form Wald test. We document~\eqref{eq:fs} here because the F-statistic
from~\eqref{eq:fs} is the relevant power diagnostic, and because the pattern
of $\hat\pi^{\ell}_j$ across sectors is the natural mechanism narrative.
\end{remark}

\paragraph{Why the cross-office form solves the FE-absorption problem.}
The standalone tier shifts $\Align^{\mathrm{G}}$ and $\Align^{\mathrm{P}}$
have weak cross-municipality variation that survives the standard FE
structure $\gamma_{jm}+\alpha_{jt}$ in~\eqref{eq:fs}: $\Align^{\mathrm{P}}_{pt}$
is constant across $m$ at fixed $t$ (a national indicator), and
$\Align^{\mathrm{G}}_{s(m),pt}$ is constant within each (state,$t$) cell.
In a tier-stacked instrument $Z^{\mathrm{P}}_{jmt}$ or
$Z^{\mathrm{G}}_{jmt}$, all surviving identifying variation must come from
heterogeneity in the pre-period exposure shares $w^{\ell}_{jmp,t}$, which
are themselves persistent and largely absorbed by $\gamma_{jm}$. The
cross-office shifts in~\eqref{eq:align-MP}--\eqref{eq:align-MGP} inherit the
full cross-municipality variation of $\Align^{\mathrm{M}}$ and gate it on
the higher-tier coalition state, producing instruments with rich
cross-muni \emph{and} cross-time variation that survive any plausible FE
structure short of muni$\times$year FE (which would absorb the endogenous
variable itself). This is a mechanical reason---in addition to the
substantive one above---to prefer the cross-office instruments to the
additive tier stack.

%% =====================================================================
\section{The AR test: reduced form}\label{sec:ar}
%% =====================================================================

\subsection{Derivation}

\paragraph{Step 1: Structural equation.}
At the muni-year level, restate~\eqref{eq:struct} as
\begin{equation}\label{eq:second}
  Y_{mt}
  \;=\;
  \alpha_m + \delta_t + \beta' s_{mt}
  \;+\;
  \lambda\,\Vol_{mt}
  \;+\;
  \varepsilon_{mt}.
\end{equation}
Stack the $K$ sector-tier instruments as in~\eqref{eq:stacked}. Aggregate the
sector-level first stage~\eqref{eq:fs} to the muni-year level by reading
$s_{mt}\in\mathbb{R}^J$ as a vector and writing the muni-year first stage in
matrix form:
\begin{equation}\label{eq:fs-stacked}
  s_{mt}
  \;=\;
  \alpha^{(s)}_m + \delta^{(s)}_t
  \;+\;
  \Pi\, Z_{mt}
  \;+\;
  \Phi\, \EC_{mt}
  \;+\;
  \mu\,\Vol_{mt}
  \;+\;
  \nu_{mt},
\end{equation}
where $\Pi\in\mathbb{R}^{J\times K}$ is the matrix of first-stage coefficients
collected from~\eqref{eq:fs} after the sector-by-sector slopes have been
expanded into the wide format aligned with~\eqref{eq:stacked}, and
$\EC_{mt}\equiv(\EC^{\ell}_{jm,t})_{\ell,j}$ is the stacked exposure-control
vector.

\paragraph{Step 2: Reduced form by substitution.}
Substitute~\eqref{eq:fs-stacked} into~\eqref{eq:second} and collect terms:
\begin{align}
  Y_{mt}
  &= \alpha_m + \delta_t + \beta'\!\left(\alpha^{(s)}_m + \delta^{(s)}_t
       + \Pi\, Z_{mt} + \Phi\,\EC_{mt} + \mu\,\Vol_{mt} + \nu_{mt}\right)
       + \lambda\,\Vol_{mt} + \varepsilon_{mt} \notag \\
  &= \tilde\alpha_m + \tilde\delta_t
       + \underbrace{(\Pi'\beta)'}_{\equiv\,\gamma'\in\mathbb{R}^K}\, Z_{mt}
       + \underbrace{(\Phi'\beta)'}_{\equiv\,\psi'}\, \EC_{mt}
       + \underbrace{(\beta'\mu+\lambda)}_{\equiv\,\tilde\lambda}\,\Vol_{mt}
       + \underbrace{\beta'\nu_{mt}+\varepsilon_{mt}}_{\equiv\,\eta_{mt}},
       \label{eq:rf}
\end{align}
where $\tilde\alpha_m=\alpha_m+\beta'\alpha^{(s)}_m$ and
$\tilde\delta_t=\delta_t+\beta'\delta^{(s)}_t$ are absorbed by the
muni and year fixed effects. Equation~\eqref{eq:rf} is the
\emph{reduced-form GDP equation}.

\paragraph{Step 3: $H_0$ in the reduced form.}
Under $H_{0}:\beta=\bzero$,
\begin{equation}\label{eq:rf-null}
  \gamma \;=\; \Pi'\beta \;=\; \bzero \in \mathbb{R}^{K},
  \qquad
  \psi \;=\; \Phi'\beta \;=\; \bzero,
  \qquad
  \tilde\lambda \;=\; \lambda.
\end{equation}
Equation~\eqref{eq:rf-null} is the testable content. It says that under
$H_0$ all $K$ instruments enter the reduced-form GDP equation with zero
coefficient. This is true \emph{regardless of $\Pi$}: the implication holds
whether the first stage is strong, weak, or zero.

\paragraph{Step 4: The AR test as the joint Wald on $Z_{mt}$.}
Estimate
\begin{equation}\label{eq:rf-est}
  Y_{mt}
  \;=\;
  \tilde\alpha_m + \tilde\delta_t
  + \gamma'\,Z_{mt}
  + \psi'\,\EC_{mt}
  + \tilde\lambda\,\Vol_{mt}
  + \eta_{mt},
\end{equation}
and test
\begin{equation}\label{eq:wald}
  H_{0}^{\mathrm{AR}}:\quad \gamma = \bzero \quad\Longleftrightarrow\quad H_0:\;\beta=\bzero.
\end{equation}
The cluster-robust Wald statistic for $\gamma=\bzero$ in~\eqref{eq:rf-est} is
the AR test statistic. Operationally:
\begin{quote}
\texttt{mod\_rf <- feols(log\_gdp\_pc \textasciitilde{} Z + EC + Vol \texttt{|} muni\_id + year,\\
\hphantom{mod\_rf <- }data = dt, vcov = \textasciitilde muni\_id)}\\
\texttt{ar\_stat <- fixest::wald(mod\_rf, keep = "\textasciicircum Z\_")\$stat}
\end{quote}

\subsection{Equivalence with the canonical AR statistic}

The canonical Anderson--Rubin statistic of \citet{anderson1949estimation} for
$H_0:\beta=\beta_0$ in the simultaneous equation $Y=X\beta+u$ with instruments
$Z$ is
\begin{equation}\label{eq:canonical-ar}
  \mathrm{AR}(\beta_0)
  \;=\;
  \frac{(Y - X\beta_0)' P_Z (Y - X\beta_0)/K}
       {(Y - X\beta_0)' M_Z (Y - X\beta_0)/(N-K)},
\end{equation}
with $P_Z=Z(Z'Z)^{-1}Z'$ and $M_Z=I-P_Z$ after projection of all variables on
the fixed-effect space. Under correct specification and $\beta=\beta_0$,
$\mathrm{AR}(\beta_0)\sim F_{K,\,N-K-n_{\mathrm{FE}}}$ exactly under
homoskedasticity, and asymptotically $\chi^2_K/K$ under
heteroskedasticity-robust forms \citep{andrews2019weak,kleibergen2002pivotal}.

\begin{lemma}[Reduced-form equivalence at $\beta_0=\bzero$]\label{lem:ar-rf}
At $\beta_0=\bzero$, the AR statistic~\eqref{eq:canonical-ar} reduces to the
$F$-statistic for $\gamma=\bzero$ in the reduced-form
regression~\eqref{eq:rf-est}.
\end{lemma}
\begin{proof}[Sketch]
At $\beta_0=\bzero$, $Y-X\beta_0=Y$ and \eqref{eq:canonical-ar} becomes the
ratio of the explained sum of squares of $Y$ on $Z$ (divided by $K$) to the
residual sum of squares (divided by $N-K$), which is exactly the
$F$-statistic for the joint significance of $Z$ in
$Y=\gamma' Z + \mathrm{(FE)} + \eta$. With heteroskedasticity-robust /
cluster-robust variance, the corresponding Wald form replaces the homoskedastic
$F$ and is the cluster-robust AR statistic of
\citet{finlay2009implementing}.
\end{proof}

\paragraph{What this buys us.}
\autoref{lem:ar-rf} is what makes the AR optimality test cheap to implement:
the practitioner only needs to estimate~\eqref{eq:rf-est} and run a Wald
test on the instruments. No first stage is needed for the test; no $\hat\beta$
needs to exist. The first stage~\eqref{eq:fs}--\eqref{eq:fs-stacked} is
informative for power but not for size.

\subsection{Weak-instrument robustness, many instruments, and CLR}

The AR test has correct size under $H_0:\beta=\bzero$ for any value of $\Pi$,
including $\Pi=0$. This follows from the fact that under $H_0$ the
reduced-form coefficient $\gamma=\Pi'\beta$ vanishes regardless of $\Pi$. The
2SLS Wald test is not size-correct when $\Pi$ is small, because its
critical values rely on a non-degenerate first-stage limit
\citep{andrews2019weak,dufour1997some}.

For $K>1$, the conditional likelihood ratio (CLR) test of
\citet{moreira2003conditional} dominates AR in power for many parameter
configurations. We treat CLR as a complementary robustness statistic; it is
not yet implemented in the project pipeline.

For $K\to\infty$ asymptotics (the secondary margin
$\texttt{cnae\_section}\times\texttt{S3}$ with $K\approx 50$--$150$),
\citet{mikusheva2022many} provide the ridge-regularised jackknifed AR test that
is valid for $K/N\to c\in(0,1)$. This is the relevant inference for the
many-instrument robustness specification.

%% =====================================================================
\section{Volume control}\label{sec:vol}
%% =====================================================================

The volume term $\Vol_{mt}$ in~\eqref{eq:struct} partitions the political
shock's effect on $Y_{mt}$ into a level (or volume) channel and a composition
channel. Without it, the test of $H_0:\beta=\bzero$ would conflate the two:
turnover both reallocates credit across sectors and changes its aggregate
amount, and~\eqref{eq:dirderiv} concerns only the former.

\paragraph{Construction.}
\begin{equation}\label{eq:vol-explicit}
  \Vol_{mt} \;\equiv\;
  \frac{\sum_{j}\sum_{f\in\cF(j,m)}\mathrm{BNDES}_{fmt}}{\mathrm{GDP}_{m,2002}}
  \;=\;
  \frac{\mathrm{BNDES\_total}_{mt}}{\mathrm{GDP}_{m,2002}}.
\end{equation}
The numerator is current-period total municipal BNDES disbursements; the
denominator is initial-period municipal GDP (2002 R\$). The ratio is unit-free
and time-comparable across municipalities. The choice of the 2002 denominator
fixes the scaling at a pre-treatment value, so $\Vol_{mt}$ varies in $t$ only
through the numerator.

\begin{remark}[Why a pre-treatment denominator]
Using contemporaneous $\mathrm{GDP}_{mt}$ as the denominator would induce
$\mathrm{cor}(\Vol_{mt},Y_{mt})\neq 0$ mechanically through the
denominator, bias the partialling-out of the volume channel, and create a
``bad-control'' problem for the AR test. The 2002 denominator avoids this.
\end{remark}

\paragraph{Should $\Vol_{mt}$ be instrumented?}
$\Vol_{mt}$ is itself endogenous: the same political shocks that move sector
shares move the aggregate. The natural instrument for $\Vol_{mt}$ is a
muni-level aggregation of $Z^{\ell}_{jmt}$, weighted by the pre-period sector
shares; this is the muni-level shift-share for the level channel. The four
candidate operationalisations from the strategy memo (A10) are:
\begin{enumerate}[leftmargin=*,nosep]
\item \emph{Pure AR.} OLS for both $s_{mt}$ and $\Vol_{mt}$ in the reduced form
  (i.e., neither is treated as the structural endogenous variable). The AR test
  is on $Z_{mt}$; $\Vol_{mt}$ is included as a pre-determined control.
\item \emph{Partial IV} (\textbf{baseline}, per D24). Treat $s_{mt}$ as
  endogenous and instrumented by $Z_{mt}$; treat $\Vol_{mt}$ as exogenous.
  The reduced-form AR equation~\eqref{eq:rf-est} includes $\Vol_{mt}$ on the
  right-hand side and tests $\gamma=\bzero$ on $Z_{mt}$.
\item \emph{Full IV.} Instrument both $s_{mt}$ and $\Vol_{mt}$. This requires a
  separate instrument for $\Vol_{mt}$ that is not collinear with $Z_{mt}$ in
  finite sample; the obvious candidate is a sum or PCA of $Z^{\ell}_{jmt}$
  across $j$, but that creates near-collinearity with $Z_{mt}$ itself.
\item \emph{Mixed.} OLS on $s_{mt}$, IV on $\Vol_{mt}$. Symmetric to (3) and
  similarly suffers from near-collinearity.
\end{enumerate}
The baseline (2) is the operative specification of the
paper. Approaches (1), (3), (4) enter as robustness. The choice between (1)
and (2) is consequential only if $\Vol_{mt}$ is a strong correlate of
$\varepsilon_{mt}$ \emph{after} partialling out $Z_{mt}$ and the FE; under the
exclusion restriction this residual correlation is bounded by the strength of
non-political confounders of total credit.

\begin{remark}[The collinearity tradeoff]
Approaches (3) and (4) deliver consistent estimators of both $\beta$ and
$\lambda$ if a separate instrument for $\Vol_{mt}$ exists. The candidate
$\sum_{j}Z^{\ell}_{jmt}$ collapses across sectors and discards the
identifying variation in cross-sector tilts that is precisely what powers the
AR test. Building a non-trivial second instrument requires new variation
(e.g.\ a national fiscal-cycle shock that scales the BNDES envelope without
shifting its sectoral allocation). We do not pursue this here.
\end{remark}

%% =====================================================================
\section{Controls and fixed effects}\label{sec:fe}
%% =====================================================================

We document the role of each control / fixed effect and the assumption it
enables.

\paragraph{Municipality fixed effect $\alpha_m$.}
Absorbs all time-invariant municipal characteristics: geography, fixed
infrastructure stocks, time-invariant institutional quality, fixed
sectoral-comparative-advantage components of $\theta_{m}$. With $\alpha_m$,
identification of $\beta$ uses only \emph{within}-muni variation in
$s_{mt}$. Required for \autoref{cond:BJS1} to be plausible: a level
correlation between alignment frequency and economic level across
municipalities is allowed.

\paragraph{Year fixed effect $\delta_t$.}
Absorbs national-level shocks affecting all municipalities in the same year:
common monetary policy, national IPCA inflation, national electoral
cycles affecting all jurisdictions simultaneously. Required so the
identifying variation in $\Align^{\ell}_{mpt}$ is the cross-muni dispersion in
who-rules-where in a given year, not the national alternation between Lula,
Dilma, and Temer.

\paragraph{Exposure control $\EC^{\ell}_{jm,t}$.}
$\EC^{\ell}_{jm,t}=\sum_p w^{\ell}_{jmp,t}$ measures total pre-period political
connectedness of $(j,m)$ at office $\ell$. It is the level-restricted analogue
of \citeauthor{goldsmith2020bartik}'s share-controls strategy: it absorbs any
direct effect of the share component of the SSIV that is not mediated by the
shift component. In the AR reduced form~\eqref{eq:rf-est}, $\EC^{\ell}_{jm,t}$
is included alongside $Z^{\ell}_{jmt}$ to remove this share-channel
confounding from the test on $\gamma$.

\paragraph{Pre-period GDP trend (robustness).}
A muni-specific linear trend in $t$, fitted on $Y_{m,t-1}$ and earlier, can be
included to harden \autoref{cond:BJS1} against persistent muni-level growth
differentials that correlate with party alignment. This is the practical
manifestation of the $q_k$ in \autoref{cond:BJS1}. We treat it as robustness
because it costs degrees of freedom (1 parameter per municipality) and
because the muni-by-year FE design already absorbs everything cross-cycle.

\paragraph{State$\times$year fixed effect (robustness).}
$\zeta_{s(m)t}$ where $s(m)$ is the state of $m$ absorbs governor-level shocks
common to all municipalities in a state. Including it removes the variation
exploited by the governor-tier instruments $Z^{\mathrm{G}}_{jmt}$, so it is
incompatible with using all three tiers; it is included as a within-state
robustness exercise that uses only mayor-tier instruments.

\paragraph{What is not included by design.}
We do not include lagged sector shares $s_{m,t-1}$ as controls. Doing so would
absorb the pre-period mix and break the within-simplex projection: the
gradient $\nabla_s Y$ is identified from changes \emph{around} the
pre-period mix, not from levels conditional on it. This is the standard
``Nickell bias''-style argument adapted to the simplex.

%% =====================================================================
\section{Inference}\label{sec:inference}
%% =====================================================================

\paragraph{Primary clustering.}
Standard errors in~\eqref{eq:rf-est} are clustered by $\texttt{muni\_id}$.
The number of clusters ($M=5{,}570$) is large relative to the parameter
dimension and to the number of effective shocks per cycle, so cluster-robust
inference is well-calibrated.
\citet{finlay2009implementing} establish that the cluster-robust Wald
statistic for $\gamma=\bzero$ in~\eqref{eq:rf-est} is the cluster-robust
AR test, with size correctly controlled under quite general within-cluster
dependence structures.

\paragraph{Two-way clustering.}
For the sector-level first stage~\eqref{eq:fs} we cluster two-way by
municipality and sector. For the muni-year reduced form~\eqref{eq:rf-est}
the natural second clustering dimension is harder: clustering by
\emph{shock} ($\ell,p,m,t$) is collinear with the by-muni clustering, and
clustering by year would yield only $T=16$ clusters.

\paragraph{Spatial correlation.}
\citet{adao2019shift} show that residuals in shift-share designs are
correlated across units that share similar exposure structures (here, similar
$w^{\ell}_{jmp,t}$ profiles). Cluster-robust SEs by $\texttt{muni\_id}$ do not
fully internalize this correlation, and may understate uncertainty. Two
remedies:
\begin{enumerate}[leftmargin=*,nosep]
\item State-level clustering as a robustness check. The $27$ Brazilian states
  are coarse and trigger few-clusters concerns
  \citep{carter2017asymptotic}; we use the wild cluster bootstrap
  in this case.
\item The Adão--Kolesár--Morales correction \citep{adao2019shift}, computed at
  the shock level. This is the right inference under the BJS framework but is
  not yet implemented in the pipeline.
\end{enumerate}

\paragraph{Concern: state-level correlation and the cross-office channels.}
The standalone tier instruments $Z^{\mathrm{G}}_{jmt}$ and
$Z^{\mathrm{P}}_{jmt}$ have weak cross-municipality variation that
survives the standard FE structure---$\Align^{\mathrm{P}}_{pt}$ is
constant across municipalities at any $t$ and so is absorbed by
$\delta_t$, and $\Align^{\mathrm{G}}_{s(m),pt}$ is constant within each
(state,$t$) cell---which is the mechanical reason we drop them from the
primary stack in favor of the cross-office shifts of
\autoref{subsec:cross-office}. The cross-office instruments
$Z^{\mathrm{M{\cdot}P}}_{jmt}$, $Z^{\mathrm{M{\cdot}G}}_{jmt}$, and
$Z^{\mathrm{M{\cdot}G{\cdot}P}}_{jmt}$ inherit the cross-municipality
variation of $\Align^{\mathrm{M}}_{mpt}$ and gate it on the higher-tier
coalition state, so they survive any plausible FE structure short of
muni$\times$year FE. The Adão--Kolesár--Morales correction
\citep{adao2019shift} remains the right inference under the BJS
framework when within-state residual correlation is a concern; it is
applied at the shock level, where shocks now correspond to
channel-party-municipality-time cells $k=(c,p,m,t)$. We report
\emph{channel-specific} AR statistics for each $c\in\mathcal{C}$ alongside
the joint test, with the pure-mayoral statistic as the conservative
benchmark and the cross-office statistics as the mechanism-relevant
hypotheses.

\paragraph{Effective sample size and time-series content.}
The levels instrument $Z^{c}_{jmt}$ is a step function in $t$ that is
piecewise constant between consecutive elements of $E_c$. With muni and
year FE, the variation exploited is the cross-municipality dispersion in
instrument step-changes at the inauguration boundaries: for the pure
mayoral channel, mayoral inaugurations $\{2005,2009,2013,2017\}$; for the
cross-office channels, the union of mayoral and gubernatorial/presidential
inaugurations $\{2003,2005,2007,\ldots,2017\}$. Each municipality
contributes $\sim 3$--$4$ independent step-changes for the pure mayoral
channel and $\sim 7$--$8$ for each cross-office channel; the
cross-municipality average has $\sim M\times 3$ degrees of freedom for
the mayoral channel and $\sim M\times 7$ for each cross-office channel.
This is the relevant denominator for power, not the panel-level
$M\times T$.

%% =====================================================================
\section{Summary}\label{sec:summary}
%% =====================================================================

The Anderson--Rubin optimality test of this paper is operationally a Wald test
on the instruments in the reduced-form GDP regression~\eqref{eq:rf-est}. The
substantive content is the directional-derivative condition~\eqref{eq:dirderiv}
of \autoref{sec:theory}, which equates the local-optimality null to
$\beta=\bzero$ in the structural equation~\eqref{eq:struct}. The reduced-form
equivalence (\autoref{lem:ar-rf}) is what makes the test (i) weak-instrument
robust, (ii) joint across $K$ sector-tier instruments, and (iii) implementable
without producing a $\hat\beta$ point estimate when the first stage is weak.
The shift-share framework of \citet{borusyak2022quasi} provides the
identification conditions; \autoref{prop:lvlchg} states the formal sufficient
conditions for the levels and changes timing variants and shows that the
changes variant carries a milder requirement on pre-period dynamics. The
volume control $\Vol_{mt}$ partials out the level channel of the political
shock, leaving only the composition channel as the object of the test.
Inference is cluster-robust by municipality, with state-level and
\citet{adao2019shift} corrections in robustness.

%% =====================================================================
\appendix
\section*{Appendix: Notation glossary}

\begin{center}
\renewcommand{\arraystretch}{1.2}
\begin{tabular}{ll}
\toprule
Symbol & Object \\
\midrule
$m,j,t,\ell,p$ & municipality, sector, year, office, party \\
$f$ & firm \\
$\cF(j,m),\,\cF^{\mathrm{pre}}_{jm,e}$ & firm support of $(j,m)$ / pre-period support at inauguration $e$ \\
$s_{jmt}$ & sector $j$ employment share in $(m,t)$, $\sum_j s_{jmt}=1$ \\
$s^{\mathrm{B}}_{jmt}$ & sector $j$ BNDES credit share in $(m,t)$ \\
$Y_{mt}$ & log real GDP in $(m,t)$ \\
$\beta$ & $J$-vector of GDP gradient w.r.t.\ employment shares \\
$\Vol_{mt}$ & $\mathrm{BNDES\_total}_{mt}/\mathrm{GDP}_{m,2002}$ \\
$T^{\ell}_{t}$ & pre-election window for office $\ell$ at year $t$, eq.\eqref{eq:window} \\
$L_{f,p,t},\,L_{f,t}$ & owners of firm $f$ affiliated with party $p$ / total owners at $t$ \\
$\omega^{\ell}_{fp,t}$ & firm-level baseline party exposure, pooled-count, eq.\eqref{eq:omega-firm} \\
$N^{\ell}_{jm,t}$ & owner-count normalizer in cell $(j,m)$, eq.\eqref{eq:N-norm} \\
$w^{\ell}_{jmp,t}$ & sector-level baseline party exposure, eq.\eqref{eq:wjmp} \\
$\FA^{c}_{fmt}$ & firm-level cross-office instrument for channel $c$, eq.\eqref{eq:FA-firm} \\
$\Align^{\ell}_{mpt}$ & $\ind\{\text{party } p \text{ holds office } \ell \text{ in } m \text{ at } t\}$ \\
$\Delta\Align^{\ell}_{mpt}$ & alignment turnover ($\neq 0$ at inauguration) \\
$\Align^{c}_{mpt}$ & cross-office shift for channel $c$, eqs.\eqref{eq:align-MP}--\eqref{eq:align-MGP} \\
$Z^{c}_{jmt}$ & sector-channel levels instrument, eqs.\eqref{eq:Z-M}--\eqref{eq:Z-MGP} \\
$\Delta Z^{c}_{jmt}$ & sector-channel changes instrument, eq.\eqref{eq:dZ} \\
$Z_{mt}$ & stacked muni-year instrument vector, eq.\eqref{eq:stacked}, dim.\ $K$ \\
$\EC^{\ell}_{jm,t}$ & sector-level exposure control, $\sum_p w^{\ell}_{jmp,t}$ \\
$\Pi$ & first-stage matrix in eq.\eqref{eq:fs-stacked} \\
$\gamma=\Pi'\beta$ & reduced-form coefficient on $Z_{mt}$ in eq.\eqref{eq:rf-est} \\
$\mathrm{AR},\,F_{\mathrm{AR}}$ & Anderson--Rubin statistic / $F$-form \\
$H_0:\beta=\bzero$ & local-optimality null \\
\bottomrule
\end{tabular}
\end{center}

%% =====================================================================
\begin{thebibliography}{99}

\bibitem[Adão, Kolesár, and Morales(2019)]{adao2019shift}
Ad\~ao, R., M.~Koles\'ar, and E.~Morales (2019).
``Shift-share designs: Theory and inference,''
\emph{Quarterly Journal of Economics} 134(4): 1949--2010.

\bibitem[Anderson and Rubin(1949)]{anderson1949estimation}
Anderson, T.~W., and H.~Rubin (1949).
``Estimation of the parameters of a single equation in a complete system of
stochastic equations,''
\emph{Annals of Mathematical Statistics} 20(1): 46--63.

\bibitem[Andrews, Stock, and Sun(2019)]{andrews2019weak}
Andrews, I., J.~H. Stock, and L.~Sun (2019).
``Weak instruments in instrumental variables regression: Theory and practice,''
\emph{Annual Review of Economics} 11: 727--753.

\bibitem[Borusyak, Hull, and Jaravel(2022)]{borusyak2022quasi}
Borusyak, K., P.~Hull, and X.~Jaravel (2022).
``Quasi-experimental shift-share research designs,''
\emph{Review of Economic Studies} 89(1): 181--213.

\bibitem[Borusyak, Hull, and Jaravel(2025)]{borusyak2025practical}
Borusyak, K., P.~Hull, and X.~Jaravel (2025).
``A practical guide to shift-share instruments,''
\emph{Journal of Economic Perspectives} 39(1): 181--204.

\bibitem[Carter, Schnepel, and Steigerwald(2017)]{carter2017asymptotic}
Carter, A.~V., K.~T. Schnepel, and D.~G. Steigerwald (2017).
``Asymptotic behavior of a $t$-test robust to cluster heterogeneity,''
\emph{Review of Economics and Statistics} 99(4): 698--709.

\bibitem[Dufour(1997)]{dufour1997some}
Dufour, J.-M. (1997).
``Some impossibility theorems in econometrics with applications to structural
and dynamic models,''
\emph{Econometrica} 65(6): 1365--1387.

\bibitem[Finlay and Magnusson(2009)]{finlay2009implementing}
Finlay, K., and L.~M. Magnusson (2009).
``Implementing weak-instrument robust tests for a general class of
instrumental-variables models,''
\emph{Stata Journal} 9(3): 398--421.

\bibitem[Goldsmith-Pinkham, Sorkin, and Swift(2020)]{goldsmith2020bartik}
Goldsmith-Pinkham, P., I.~Sorkin, and H.~Swift (2020).
``Bartik instruments: What, when, why, and how,''
\emph{American Economic Review} 110(8): 2586--2624.

\bibitem[Guggenberger and Kleibergen(2012)]{guggenberger2012generalized}
Guggenberger, P., and F.~Kleibergen (2012).
``A more powerful subvector Anderson--Rubin test in linear instrumental
variables regression,''
working paper.

\bibitem[Kleibergen(2002)]{kleibergen2002pivotal}
Kleibergen, F. (2002).
``Pivotal statistics for testing structural parameters in instrumental
variables regression,''
\emph{Econometrica} 70(5): 1781--1803.

\bibitem[Mikusheva and Sun(2022)]{mikusheva2022many}
Mikusheva, A., and L.~Sun (2022).
``Inference with many weak instruments,''
\emph{Review of Economic Studies} 89(5): 2663--2686.

\bibitem[Moreira(2003)]{moreira2003conditional}
Moreira, M.~J. (2003).
``A conditional likelihood ratio test for structural models,''
\emph{Econometrica} 71(4): 1027--1048.

\bibitem[Samuelson(1947)]{samuelson1947foundations}
Samuelson, P.~A. (1947).
\emph{Foundations of Economic Analysis}, Cambridge: Harvard University Press.

\end{thebibliography}

\end{document}
